% Solutions to the Chapter 1 exercises that are not code, from lectures/L01/appendix/c_exercises.md.
% Exercises 1.6 and 1.7 are Code and are not answered; see chapter.tex for why.

\section{Chapter 1: Embedded Linux, and the Licence You Ship With It}
\label{sol:sec:c1}
Answers to the exercises in \secref{c1:sec:exercises}. Exercises~\ref{c1:ex:6} and~\ref{c1:ex:7}
are Code, and are checked by running them: the first against the total its Check yourself line
gives, the second on the target.

% ------------------------------------------------------------------------------------------
\solution[fillaccent2]{c1:ex:1}{The four pieces}{Recall}

The \textbf{bootloader} brings up DRAM, finds the kernel, and hands it a device tree. The
\textbf{kernel} drives the hardware and provides the system call interface. The \textbf{root
filesystem} is everything that runs as a process, starting with PID 1. The \textbf{toolchain}, a
compiler, a C library, a linker and a set of headers, builds the other three and never ships, and
it fixes the ABI they share.

\xpart{a} The kernel. It is stored in the root filesystem, and that is the plausible wrong answer:
it is built by the kernel's build from the kernel's tree, it runs in the kernel's address space,
and it loads only into the kernel it was built for, which is why its path carries the version,
6.12.30. Where a file is stored does not decide which piece it belongs to.

\xpart{b} The bootloader: \file{u-boot.img} is U-Boot itself, step 3 of the chain in
\secref{c1:app:a2}.

\xpart{c} The toolchain. It lives on the build machine, and nothing on the device ever runs it.

\xpart{d} The root filesystem. \file{/init} is PID 1, the first userspace process; on this target it
is the shell script \file{tools/target/init}, which \file{ci/rootfs.sh} copies in.

\xpart{e} The kernel, although it travels beside the kernel rather than inside it. Its source lives
in the kernel tree, the kernel's build compiles it, and it is written against the bindings the
kernel documents, so it changes when the kernel does. The bootloader is what loads it and hands it
over (\secref{c1:app:a3}), and on a real board U-Boot often edits it on the way, filling in the
memory size and the command line, so ``the bootloader'' is a defensible answer if it is argued that
way. On this target it is neither: QEMU writes it (\secref{c1:app:a4}).

\xpart{f} The root filesystem. \file{/bin/busybox} is the whole userland of this target.

% ------------------------------------------------------------------------------------------
\solution[fillaccent2]{c1:ex:2}{Why there are two bootloaders}{Recall}

\xpart{a} U-Boot is around a megabyte and must run from DRAM, and DRAM does not work until its
controller has been configured, while the only memory there is at power-on is a few tens of
kilobytes of SRAM on the SoC. The SPL is a bootloader small enough to run from that SRAM, and its
whole job is to initialise the DRAM controller and load U-Boot into the DRAM it has just brought up.

\xpart{b} The boot stops at the first handover, from the boot ROM to the next stage, before U-Boot
has printed a line. The ROM loads the next stage into SRAM from wherever it is told to look; with
the SPL deleted it finds either nothing it accepts or U-Boot, which does not fit in SRAM and needs
the DRAM nobody has initialised. The board produces no bootloader output at all and looks dead, and
the space saved is the SPL's, around 64 KB.

\xpart{c} Recoverable: a boot ROM that falls back to loading over USB or a UART, or from a removable
SD card, when the flash holds nothing it can boot, or strap pins that force it to; a bad bootloader
can then always be replaced from outside. Not recoverable: a ROM that boots only from soldered flash
and has no fallback. One bad write to the bootloader partition and the board is a coaster, short of
a JTAG probe or a soldering iron.

% ------------------------------------------------------------------------------------------
\solution[fillaccent3]{c1:ex:3}{What the toolchain triplet promises}{Hand calculation}

\xpart{a} \code{aarch64} is the architecture: 64-bit ARM, little-endian, since the big-endian
variant is \code{aarch64_be}. \code{linux} is the operating system the output runs on: it makes
Linux system calls and follows Linux's conventions for an executable. \code{gnu} is the C library
and the ABI: glibc, and the GNU/Linux ABI for AArch64, whose dynamically linked programs are started
by \file{/lib/ld-linux-aarch64.so.1}. The full form has a vendor field as well,
\code{aarch64-unknown-linux-gnu}, which the short name leaves out.

\xpart{b} \code{arm} is 32-bit ARM, in practice a Cortex-M or Cortex-R microcontroller. \code{none}
is the operating system, and there is none: the binary is written for bare metal. \code{eabi} is the
ARM Embedded ABI, the calling convention and object format for code with no operating system beneath
it; no C library is named, and in practice it is newlib. So it is meant for a microcontroller, or
for a stage that runs before any operating system does. It cannot run on Linux as a program: nothing
in it makes a Linux system call, its startup code and linker script assume it owns the machine from
a fixed address, and there is no Linux loader or C library for it to be started by.

\xpart{c} It builds, and fails at run time, when the target's dynamic loader starts it. glibc
versions its symbols, and a program records the versions of every symbol it was linked against;
since glibc 2.34 even \code{__libc_start_main}, which every dynamically linked program calls,
carries the version \code{GLIBC_2.34}. The 2.31 on the target does not have it, so the loader
refuses before \code{main} runs, with a message of the form \code{version `GLIBC_2.34' not found
(required by ./prog)}. Nothing at build time can see this coming, because the build only ever saw
2.36. The fix is to build against a C library no newer than the oldest target's, or to link
statically, which is what this course does for its own test programs, because its target has no C
library at all.

\xpart{d} It installs the package for the build machine's architecture: its headers in
\file{/usr/include}, and an x86-64 library in \file{/usr/lib/x86_64-linux-gnu}. Debian's
cross-compiler searches \file{/usr/include} (\secref{c1:app:a5}), so the missing header may even be
found, and that is the trap: the failure moves to the link, which cannot find \code{-lfoo}, because
the only \code{libfoo} is an x86-64 library in a directory the cross linker never searches and could
not link against if it did. If the missing header is one of the few that differ by architecture,
which Debian keeps in \file{/usr/include/x86_64-linux-gnu}, the error does not move at all. What
works is to install the library \emph{for the target}: on Debian, \code{dpkg --add-architecture
arm64}, an \code{apt update}, and then \code{apt install libfoo-dev:arm64}, which puts the arm64
library in \file{/usr/lib/aarch64-linux-gnu}, where the cross-compiler does look; or cross-compile
\code{libfoo} and install it into the toolchain's sysroot. For this course's target, whose root
filesystem holds no libraries, the program must also link it statically.

% ------------------------------------------------------------------------------------------
\solution[fillaccent4]{c1:ex:4}{What leaves the building}{Design}

\xpart{a} Row by row:
\begin{itemize}
  \item \textbf{U-Boot}: its complete corresponding source, meaning the version you shipped with
    your board port in it, the configuration you built it with, and the scripts and instructions
    that build it. It is ``or later'', so you may comply on GPL-2.0's terms.
  \item \textbf{The kernel}: the same, at the exact version shipped with every patch applied,
    including the driver you added, which is in the tree and therefore part of the work; the
    \code{.config}; the device tree sources; and which toolchain and which commands build it
    (\secref{c1:app:b5}).
  \item \textbf{BusyBox}: its source at the version shipped, and the \code{.config} you built it
    with, which is one of the scripts that control its compilation.
  \item \textbf{\code{libcurl}}: no source. Its licence is permissive and asks only that its
    copyright and permission notice go with every copy, which for a binary means your
    documentation.
  \item \textbf{The logging daemon}: nothing. It is a separate process making system calls
    (\secref{c1:app:b3}).
  \item \textbf{Your kernel module}: that depends on the choice, and the choice has consequences.
    Declared GPL, it is published like the kernel. Declared proprietary, it is published nowhere,
    and you carry what \secref{c1:app:b7} lists, a tainted kernel, no \code{EXPORT_SYMBOL_GPL}
    symbols and its maintenance against a kernel API that changes every release, and a
    derivative-work argument that is contested rather than settled in your favour;
    Exercise~\ref{c1:ex:5} works that argument through.
\end{itemize}

\xpart{b} No. The obligation attaches to distributing the binary, not to modifying it: whoever
receives BusyBox from you is entitled to its source whether or not you changed a line. Being
unmodified makes compliance easy, since the answer is the upstream release at the version you
shipped and your \code{.config}; it does not make it optional.

\xpart{c} Neither does. \code{libcurl}'s licence is permissive, so linking to it, dynamically or
statically, puts no obligation on the daemon's source, only the notice from a). The \code{ioctl} is
a system call: the daemon uses your module through the kernel's system call interface, which is
exactly the boundary the syscall note draws, and \secref{c1:app:b3} names \code{ioctl} in its list.
The module is yours in any case, so nobody else's licence reaches the daemon through it.

\xpart{d} Keep, for every release you shipped, the complete corresponding source of every GPL
component exactly as it was built: the kernel with its patches, your driver, its \code{.config} and
its device tree sources; U-Boot with the board port and its configuration; BusyBox and its
\code{.config}; your module, if you made it GPL; and the build scripts and the toolchain, or enough
to reproduce them, that turned all of that into the image you shipped. If you satisfy the GPL with a
written offer, GPL-2.0 section~3(b) makes the offer valid for at least three years and open to any
third party, not only to customers, and every unit you ship carries it, so in practice you keep the
source until three years after the last unit ships (\secref{c1:app:b5}). The engineering consequence
is that your build must still work then: the toolchain archived rather than named, every source
pinned, the \code{.config} kept as a build artefact, and a record of which source went into which
shipped unit.

% ------------------------------------------------------------------------------------------
\solution[fillaccent4]{c1:ex:5}{Where the boundary falls}{Design}

\xpart{a} Two works. The syscall note settles it: user programs that use kernel services by normal
system calls do not fall under the heading of derived work (\secref{c1:app:b3}). It is the
copyright holders' own statement rather than a reading of the GPL, and it is not contested.

\xpart{b} One work. Static linking puts your code and the library into one binary, and that is
broadly accepted as a derivative work of both (\secref{c1:app:b2}). Distributing it means
distributing the whole under GPL-2.0, source included.

\xpart{c} Two works, and the licence says so. The LGPL-2.1 treats a program that only links to the
library as a work that uses the library, which may carry any licence you like, provided you publish
the library's source and let the recipient relink your program against a modified copy of the
library, which a dynamic link allows as it stands. This is the question the LGPL exists to settle.

\xpart{d} Contested. The Free Software Foundation holds that dynamically linking against a GPL
library creates a derivative work; others disagree, and nothing has settled it
(\secref{c1:app:b2}). The engineering answer is not to ship it without a lawyer's opinion, or to
find a library under the LGPL or a permissive licence that does the same job.

\xpart{e} Contested, with most of the weight on one work. A module is compiled against the kernel's
headers, runs in its address space and calls its functions directly, and the widely held view among
kernel developers is that it is a derivative work (\secref{c1:app:b4}). Whatever a court would
decide, the kernel enforces a version of the argument itself: a proprietary module taints the kernel
and cannot use \code{EXPORT_SYMBOL_GPL} symbols.

\xpart{f} Two works: mere aggregation (\secref{c1:app:b2}). Sharing a filesystem image does not make
two programs one work, and nobody seriously argues otherwise.

% ------------------------------------------------------------------------------------------
\solution[fillaccent]{c1:ex:8}{An audit by hand against an audit by script}{Cross-check}

\xpart{a} Eight of the ten carry a tag, each on its first line in a \code{//} comment. Two carry
none, and their licence is in the header comment instead.

\begin{narrowtable}{@{}lll@{}}
  \thead{File} & \thead{Identifier} & \thead{Where} \\
  \midrule
  \file{rtc-ds1307.c} & \code{GPL-2.0-only} & Tag, line 1 \\
  \file{rtc-pl031.c} & \code{GPL-2.0-or-later} & Tag, line 1 \\
  \file{rtc-starfire.c} & None & Header comment: ``License: GPL'' \\
  \file{rtc-rs5c313.c} & None & Header comment: the GPL, and see \file{COPYING} \\
  \file{rtc-cmos.c} & \code{GPL-2.0-or-later} & Tag, line 1 \\
  \file{rtc-m41t80.c} & \code{GPL-2.0-only} & Tag, line 1 \\
  \file{rtc-pcf8563.c} & \code{GPL-2.0-only} & Tag, line 1 \\
  \file{rtc-abx80x.c} & \code{GPL-2.0} & Tag, line 1 \\
  \file{rtc-max77686.c} & \code{GPL-2.0+} & Tag, line 1 \\
  \file{rtc-s35390a.c} & \code{GPL-2.0-or-later} & Tag, line 1 \\
\end{narrowtable}

\noindent \code{GPL-2.0} is the old spelling of \code{GPL-2.0-only} and \code{GPL-2.0+} the old
spelling of \code{GPL-2.0-or-later}, so the ten are four GPL-2.0-only, four GPL-2.0-or-later and two
untagged: four strings, and two licences.

\xpart{b} From the sample alone: two licences, GPL-2.0-only and GPL-2.0-or-later, in about equal
shares, and about one file in five with no tag.

\xpart{c} A script that meets Exercise~\ref{c1:ex:6} reports this over the 184 files:

\begin{narrowtable}{@{}lrl@{}}
  \thead{Identifier} & \thead{Files} & \thead{Of which merged} \\
  \midrule
  \code{GPL-2.0-only} & 122 & 49 spelled \code{GPL-2.0} \\
  \code{GPL-2.0-or-later} & 59 & 28 spelled \code{GPL-2.0+} \\
  \code{LGPL-2.1+} & 1 & \\
  Untagged & 2 & \file{rtc-rs5c313.c}, \file{rtc-starfire.c} \\
  \midrule
  Total & 184 & \\
\end{narrowtable}

\noindent The counts check against the total: $122 + 59 + 1 + 2 = 184$. No file in the directory
carries a second tag, so taking the first loses nothing here. The single \code{LGPL-2.1+} is
\file{lib_test.c}, the directory's KUnit test; it is the old spelling of LGPL-2.1-or-later, which
the exercise does not require the script to merge and a careful script merges anyway.

\xpart{d}
\begin{itemize}
  \item \textbf{Strings and licences.} Five distinct strings: \code{GPL-2.0-only}, \code{GPL-2.0},
    \code{GPL-2.0-or-later}, \code{GPL-2.0+} and \code{LGPL-2.1+}. Three distinct licences:
    GPL-2.0-only, GPL-2.0-or-later and LGPL-2.1-or-later. The strings count spellings, and the
    difference is the number of licences that appear under two spellings, here the two GPL ones, so
    $5 = 3 + 2$. The string count can never be lower than the licence count, and a report that
    calls strings licences overstates what the tree is under.
  \item \textbf{The untagged files} are \file{rtc-rs5c313.c} and \file{rtc-starfire.c}, both of them
    in the sample, and the routes differ. \file{rtc-rs5c313.c}'s header says it is subject to the
    GNU GPL and sends you to \file{COPYING}, which is GPL-2.0-only; its \code{MODULE_LICENSE("GPL")}
    agrees, but cannot tell ``only'' from ``or later'' by itself, as \file{include/linux/module.h}
    says. \file{rtc-starfire.c}'s header says only ``License: GPL'', with no version, and the file
    has no \code{MODULE_LICENSE} at all, since it can only be built in; its version comes from the
    rule in \file{Documentation/process/license-rules.rst} that the licence in \file{COPYING}
    applies to the kernel source as a whole. Both end at GPL-2.0-only, one because its header says
    so and the other because nothing says otherwise.
  \item \textbf{The sample did not predict the whole.} It had 40\% GPL-2.0-only, 40\%
    GPL-2.0-or-later and 20\% untagged. The directory is $122/184 = 66.3\%$ GPL-2.0-only,
    $59/184 = 32.1\%$ GPL-2.0-or-later, $1/184 = 0.5\%$ LGPL and $2/184 = 1.1\%$ untagged. The
    sample over-counted untagged files eighteen times over, under-counted GPL-2.0-only, and missed
    a licence altogether. That is what a sample of names you recognise does: the names you know are
    the old, widely used drivers, and old files are the ones that predate the tags and still carry
    a licence in a header comment. The bulk of the directory is newer drivers tagged
    \code{GPL-2.0-only}, which are exactly the names nobody recognises. A random ten would have
    held both untagged files with probability $\frac{10 \times 9}{184 \times 183} = 90/33{,}672$,
    about 1 in 370, so this sample was chosen, not drawn.
\end{itemize}

\xpart{e} There are three separate problems, and any two answer it:
\begin{itemize}
  \item \textbf{Nobody tags a file \code{Proprietary}.} SPDX has no such identifier, and a
    proprietary file carries no tag at all, or a \code{LicenseRef-} of its own. The absence of the
    string proves nothing; the files that could be proprietary are the untagged ones, which have to
    be read, as in d).
  \item \textbf{One directory is not the tree.} The script looked at the \code{.c} and \code{.h}
    files of \file{drivers/rtc}. Every other directory, and every file that is not C, from assembly
    and device tree sources to scripts, went unexamined.
  \item \textbf{The tree is not the product.} What ships includes things with no source in the tree
    at all: firmware that drivers load at run time, much of it under licences that are not free;
    out-of-tree and vendor modules, whose \code{MODULE_LICENSE} is in a binary no scan of sources
    will see; and everything else in the image.
\end{itemize}
